%! Zeros of Polynomials: the Rational Root Theorem — Unit 3, Lesson 3.4 — Warm-Up
\documentclass[10pt]{article}
\usepackage{algebra2-article}
\usepackage{algebra2-boxes}
\newcommand{\TallMath}[1]{$\displaystyle #1\rule[-1.4em]{0pt}{3.2em}$}

\begin{document}

\pageheader{Unit 3, Lesson 3.4}{Warm-Up}
\namedateperiod

\vspace{0.10in}

\begin{notesbox}{Getting Ready: Finishing Yesterday's Problem}
\small Lesson 3.3 gave you the tools to crack any polynomial --- \emph{once somebody told you which
number to test}. Today you build that list of numbers yourself. Two pieces of arithmetic and one
unfinished problem are all it takes.

\vspace{5pt}
\textbf{1. Divisors.} List \emph{every} integer that divides each number evenly, positives and
negatives.

\vspace{2pt}
Divisors of $6$: \blank{5.0cm} \qquad Divisors of $2$: \blank{3.0cm}

\vspace{6pt}
\textbf{2. Every fraction you can build.} Take a numerator from the first list and a denominator
from the second. Write each fraction $\dfrac{p}{q}$ in lowest terms, and cross out repeats.

\vspace{3pt}
Numerators $p$: $1, 2, 3, 6$ \qquad Denominators $q$: $1, 2$

\vspace{3pt}
All the different positive values of $\dfrac{p}{q}$: \blank{7.0cm}

\vspace{3pt}
How many values is that, counting both signs? \blank{1.4cm}

\vspace{6pt}
\textbf{3. Unfinished business.} Yesterday your group tested the divisors of $6$ on
$h(x) = 2x^3 - 3x^2 - 11x + 6$ and found that $r = -2$ and $r = 3$ both gave $0$. Divide out $(x-3)$.

\vspace{2pt}
The divisor is $x - 3$, so $r = $ \blank{1.0cm}
\begin{center}
\renewcommand{\arraystretch}{1.6}
\begin{tabular}{r|C{1.4cm}C{1.4cm}C{1.4cm}C{1.4cm}}
    & $2$ & $-3$ & $-11$ & $6$ \\
    &     &      &       &     \\ \cline{2-5}
    &     &      &       &     \\
\end{tabular}
\end{center}

\vspace{2pt}
Depressed polynomial: \blank{3.0cm} \quad Remainder: \blank{1.2cm}

\vspace{3pt}
Factor that quotient (Lesson 3.2 tools): \blank{3.4cm} \; so $h(x) = $ \blank{5.0cm}

\vspace{3pt}
The three zeros of $h$ are \blank{4.0cm}

\vspace{6pt}
\textbf{4. The gap.} Your candidate list yesterday was $\pm 1, \pm 2, \pm 3, \pm 6$ --- the divisors
of the constant term $6$. One zero of $h$ is \emph{not} on that list.

\vspace{2pt}
Which one? \blank{1.6cm} \; Where do you think its denominator came from? \blank{5.4cm}

\end{notesbox}

\end{document}
